% Qiushi Research Report Template
% Compile with LuaLaTeX.
% If the LuaLaTeX font cache is unavailable, prefix the compile command with TEXMFVAR=/tmp/qiushi_texmf_var.
% Keep figure files under figures/ or ../../workspace/figures/.
% Use refs.bib with \cite{...} when references are used.
% LaTeX supplies section numbers; section titles contain title text only.
% Long paths, URLs, and commands can use \qpath, \qurl, and \qcode.
\documentclass[11pt,a4paper]{article}

% ── Report Metadata (agent edits these only) ────────────────
\providecommand{\ReportTitle}{Report Title}
\providecommand{\ReportSubtitle}{Subtitle (optional)}
\providecommand{\ReportAuthor}{Qiushi Discovery Engine}
\providecommand{\ReportOrganization}{Qiushi Discovery Engine}
\providecommand{\ReportClassification}{Public}
\providecommand{\ReportVersion}{v1.0}
\providecommand{\ReportDate}{\today}

% ── Math + Units ────────────────────────────────────────────
\usepackage{amsmath,amssymb}
\usepackage{siunitx}
\sisetup{detect-all=true,per-mode=symbol}

% ── Graphics ────────────────────────────────────────────────
\usepackage{graphicx}
\usepackage[dvipsnames,table]{xcolor}
\usepackage{tikz}
\usepackage{pgfplots}\pgfplotsset{compat=1.18}
\graphicspath{{figures/}{../../workspace/figures/}}

% ── Tables (richer for reports) ─────────────────────────────
\usepackage{booktabs,longtable,array,tabularx,multirow,makecell,threeparttable}

% ── Bib guard ──────────────────────────────────────────────
% Reports default to plain numeric, so we use built-in \cite directly.
\newif\ifQiushiHasCitations
\let\QiushiOriginalCite\cite
\renewcommand{\cite}{\global\QiushiHasCitationstrue\QiushiOriginalCite}
\let\QiushiOriginalNocite\nocite
\renewcommand{\nocite}{\global\QiushiHasCitationstrue\QiushiOriginalNocite}

% ── Layout ──────────────────────────────────────────────────
\usepackage[inner=2.5cm,outer=2.5cm,top=2.6cm,bottom=2.7cm,headheight=22pt]{geometry}
\usepackage{setspace}\setstretch{1.15}
\setlength{\parindent}{0pt}
\setlength{\parskip}{0.6em plus 0.1em}
\raggedbottom
\emergencystretch=2em
\clubpenalty=10000
\widowpenalty=10000

\usepackage{microtype}
\usepackage{caption,subcaption}
\usepackage{enumitem}\setlist{itemsep=0.3em,topsep=0.5em,parsep=0.1em}
\usepackage{float}
\usepackage{afterpage}
\usepackage{listings}

% ── Report Palette (8 colors + Qiushi blue brand) ──────────
\definecolor{ReportPrimary}{HTML}{1F3A5F}
\definecolor{ReportAccent}{HTML}{8A5A2B}
\definecolor{ReportMuted}{HTML}{6F7782}
\definecolor{ReportLight}{HTML}{F4F7FA}
\definecolor{ReportRule}{HTML}{D6DDE5}
\definecolor{ReportFinding}{HTML}{2D5A87}
\definecolor{ReportRecommend}{HTML}{2D7A4F}
\definecolor{ReportRisk}{HTML}{B7472A}
\definecolor{ReportInsight}{HTML}{6B4E94}
% ZJU Qiushi blue (footer brand mark)
\definecolor{QiushiBlue}{HTML}{003F88}

% ── Section style (bold + color + thin rule) ───────────────
\usepackage{titlesec}
\titleformat{\section}
  {\Large\bfseries\color{ReportPrimary}}
  {\thesection}{0.6em}{}
  [\vspace{0.1em}{\color{ReportRule}\rule{\textwidth}{0.4pt}}]
\titleformat{\subsection}{\large\bfseries\color{ReportPrimary}}{\thesubsection}{0.5em}{}
\titleformat{\subsubsection}{\normalsize\bfseries}{\thesubsubsection}{0.5em}{}

% ── Footer (Nature standard: brand left, page right, rule above) ─
% No top header. Bottom: Qiushi Engine | <page>, separated by thin rule.
\usepackage{fancyhdr}
\pagestyle{fancy}
\fancyhf{}
\fancyfoot[L]{\small\color{QiushiBlue}\textbf{Qiushi Engine}}
\fancyfoot[R]{\small\thepage}
\renewcommand{\headrulewidth}{0pt}
\renewcommand{\footrulewidth}{0.4pt}
\renewcommand{\footrule}{\hbox to\headwidth{\color{ReportRule}\leaders\hrule height \footrulewidth\hfill}}
\fancypagestyle{plain}{%
  \fancyhf{}%
  \fancyfoot[L]{\small\color{QiushiBlue}\textbf{Qiushi Engine}}%
  \fancyfoot[R]{\small\thepage}%
  \renewcommand{\headrulewidth}{0pt}%
  \renewcommand{\footrulewidth}{0.4pt}%
  \renewcommand{\footrule}{\hbox to\headwidth{\color{ReportRule}\leaders\hrule height \footrulewidth\hfill}}%
}

% ── Captions ────────────────────────────────────────────────
\captionsetup{font=small,labelfont={bf,color=ReportPrimary},labelsep=quad,skip=0.5em}

% ── Banded Tables (default) ─────────────────────────────────
\rowcolors{2}{ReportLight}{white}

% ── Links ────────────────────────────────────────────────
\usepackage{xurl}
\usepackage{hyperref}
\hypersetup{colorlinks=true,linkcolor=ReportPrimary,citecolor=ReportAccent,urlcolor=ReportPrimary,
            pdftitle={\ReportTitle}}
\usepackage{bookmark}
\usepackage{cleveref}

% ── Long paths, URLs, commands, and wrapping tables ─────────
\DeclareUrlCommand\qpath{\urlstyle{tt}\small}
\DeclareUrlCommand\qurl{\urlstyle{same}\small}
\newcommand{\qcode}[1]{{\small\ttfamily\detokenize{#1}}}
\Urlmuskip=0mu plus 2mu
\newcolumntype{Y}{>{\raggedright\arraybackslash\hspace{0pt}}X}
\newcolumntype{L}[1]{>{\raggedright\arraybackslash\hspace{0pt}}p{#1}}
\lstset{
  basicstyle=\ttfamily\small,
  breaklines=true,
  breakatwhitespace=false,
  columns=fullflexible,
  keepspaces=true
}

% ── Scientific callout boxes: findings / insights / notes ──
\usepackage[most]{tcolorbox}
\tcbuselibrary{breakable,skins}
\newtcolorbox{keyfindingbox}[1][Key Finding]{
  breakable, enhanced,
  colback=ReportFinding!7, colframe=ReportFinding,
  coltitle=white, fonttitle=\bfseries,
  title=#1, boxrule=0.6pt, arc=2mm,
  left=10pt,right=10pt,top=8pt,bottom=8pt}
\newtcolorbox{insightbox}[1][Insight]{
  breakable, enhanced,
  colback=ReportInsight!7, colframe=ReportInsight,
  coltitle=white, fonttitle=\bfseries,
  title=#1, boxrule=0.6pt, arc=2mm,
  left=10pt,right=10pt,top=8pt,bottom=8pt}
\newtcolorbox{recommendbox}[1][Recommendation]{
  breakable, enhanced,
  colback=ReportRecommend!7, colframe=ReportRecommend,
  coltitle=white, fonttitle=\bfseries,
  title=#1, boxrule=0.6pt, arc=2mm,
  left=10pt,right=10pt,top=8pt,bottom=8pt}
\newtcolorbox{riskbox}[1][Note]{
  breakable, enhanced,
  colback=ReportRisk!7, colframe=ReportRisk,
  coltitle=white, fonttitle=\bfseries,
  title=#1, boxrule=0.6pt, arc=2mm,
  left=10pt,right=10pt,top=8pt,bottom=8pt}
\newcommand{\QiushiCallout}[2]{\begin{keyfindingbox}[#1]#2\end{keyfindingbox}}
% Example:
% \QiushiCallout{Key result}{State one central result or interpretation in connected prose.}
% \begin{insightbox}[Insight] Write a deep explanation or central interpretation.\end{insightbox}

% ── Cover Page ──────────────────────────────────────────────
\providecommand{\MakeReportCover}{%
  \begin{titlepage}
  \pagecolor{ReportPrimary}\afterpage{\nopagecolor}
  \color{white}
  \vspace*{2.5cm}
  \begin{flushleft}
    {\Huge\bfseries \ReportTitle\par}\vspace{0.6cm}
    {\Large\color{ReportLight}\ReportSubtitle\par}\vspace{1.5cm}
    {\color{ReportLight}\rule{6cm}{0.6pt}}
  \end{flushleft}
  \vfill
  \begin{flushleft}
    {\large \ReportAuthor\par}\vspace{0.2cm}
    {\normalsize\color{ReportLight}\ReportOrganization\par}\vspace{0.6cm}
    {\normalsize\color{ReportLight}\ReportDate\par}\vspace{0.4cm}
    {\small\color{ReportLight}Classification: \ReportClassification\quad|\quad\ReportVersion\par}
  \end{flushleft}
  \end{titlepage}}

% ============================================================
\begin{document}

\hypersetup{pageanchor=false}
\MakeReportCover
\hypersetup{pageanchor=true}
\setcounter{page}{1}

\thispagestyle{plain}
\tableofcontents
\newpage

\section*{Abstract}
\addcontentsline{toc}{section}{Abstract}

Write one continuous, self-contained summary of the research problem, core idea,
main method, decisive evidence, central result, and scientific significance. Do
not use this section as a bullet list or as a project log.

\newpage

\section{Research Background and Problem Definition}

Explain the scientific background, the open problem, why the problem matters,
and the precise research question addressed in this report. Define the main
objects, variables, physical or mathematical setting, and applicable conditions
needed for a reader to follow the rest of the report.

\section{Core Idea and Key Contributions}

Present the central insight and the real contribution of the work. Explain why
the idea is not merely a restatement of known methods, what changes in
understanding, and which results or artifacts carry the contribution.

\section{Theory and Methods}

Develop the theory, model, method, algorithm, proof structure, physical picture,
or mathematical derivation needed for the research. Keep central derivations and
scientific explanations in the main report when they are necessary for
understanding the result.

% Example equation pattern:
% \begin{equation}
%   \mathcal{L}u = f .
%   \label{eq:main-example}
% \end{equation}

\section{Experiments and Validation}

Describe the simulations, experiments, proofs, datasets, comparisons, controls,
or validation procedures that test the core idea. Explain what was measured or
computed and why it answers the research problem.

\begin{figure}[htbp]
\centering
\fbox{\rule{0pt}{1.8in}\rule{0.85\textwidth}{0pt}}
\caption{Replace this placeholder with a figure that carries a real scientific
result. Explain the figure inside the relevant scientific argument; do not
create a separate figure-by-figure catalogue.}
\label{fig:report-example}
\end{figure}

\begin{table}[htbp]
\caption{Table caption.}
\label{tab:report-example}
\begin{tabularx}{\textwidth}{YYY}
\toprule
Column A & Column B & Column C \\
\midrule
-- & -- & -- \\
-- & -- & -- \\
\bottomrule
\end{tabularx}
\end{table}

\section{Results Analysis and Discussion}

Analyze the results in connected prose. Relate figures, tables, equations, and
measurements to the scientific question. Compare alternatives, explain physical
or mathematical meaning, and state what the evidence does and does not establish.

\section{Conclusions, Significance, and Applications}

Summarize the final scientific result, why it matters, and how it may be used.
Discuss meaningful applications only when they follow from the research. Keep
the conclusion precise and grounded in the evidence.

\appendix
\section{Reproducibility Notes and Appendices}

Put concise reproducibility information, auxiliary derivations, implementation
details, source notes, or extended tables here when they are useful. Do not move
central theory, physical explanation, result analysis, or application discussion
out of the main report.

% These headings are functional roles, not a license to create many nested
% sections. The main report should be continuous scientific prose, not an
% expanded outline, blog-style Markdown, engineering log, source-note dump,
% implementation diary, or page-count filler.

\ifQiushiHasCitations
  \IfFileExists{refs.bib}{%
    \bibliographystyle{unsrt}
    \bibliography{refs}}{}
\fi

\end{document}
